\documentclass[a4paper,12pt]{article}
	\usepackage[ansinew]{inputenc}
	\usepackage{ae}
	\usepackage[brazil]{babel}
	\usepackage{lipsum}
	\usepackage{xcolor}	
	\usepackage{rotating} % <-- Permite a rotação de figuras e tabelas flutuantes
				
	\graphicspath{
		{Figuras/EPS/}
		{Figuras/JPG/}
		{Figuras/PDF/}
		{Figuras/PNG/}
		}
		
\begin{document}

	% Este arquivo ilustra a rotação (de 90 graus no sentido antihorário) de figuras e tabelas flutuantes através do pacote rotating.

	\lipsum[1-3]

	%--------------------------------------------
	% EXEMPLO 1: rotação de tabela.
	%--------------------------------------------
	\begin{sidewaystable}
		\centering
		\begin{tabular}{ccccc}
		Coluna 1 & Coluna 2 & Coluna 3 & Coluna 4 & Coluna 5 \\
		Coluna 1 & Coluna 2 & Coluna 3 & Coluna 4 & Coluna 5 \\
		Coluna 1 & Coluna 2 & Coluna 3 & Coluna 4 & Coluna 5 \\
		Coluna 1 & Coluna 2 & Coluna 3 & Coluna 4 & Coluna 5 \\
		Coluna 1 & Coluna 2 & Coluna 3 & Coluna 4 & Coluna 5 \\
		Coluna 1 & Coluna 2 & Coluna 3 & Coluna 4 & Coluna 5 \\
		\end{tabular}
		\caption{legenda da tabela.}
	\end{sidewaystable}
	
	\lipsum[4-7]
	
	%--------------------------------------------
	% EXEMPLO 2: rotação de figuras.
	%--------------------------------------------	
	\begin{sidewaysfigure}
		\centering
		\includegraphics[width=0.5\textwidth]{Tei-Gi}
		\caption{legenda da figura.}
	\end{sidewaysfigure}
	
\end{document}
